\section{Zápasenie ako šport}
\label{chapter:zapasenie}

V~tejto kapitole popisujeme zápasenie ako športovú
disciplínu so zameraním na udalosti, ktoré rozhodca
počas zápasu signalizuje a~ktoré budú predmetom
implementácie v~našej aplikácii. Postupne sa
venujeme základným pravidlám zápasu, úlohe rozhodcu,
spôsobu udeľovania bodov, napomenutiam, pasivite
a~lopatkovému víťazstvu. V~záverečnej časti
analyzujeme pravidlá z~pohľadu návrhu softvérového
riešenia.

\subsection{Definícia a~história zápasenia}

Zápasenie je kontaktný šport, ktorý kombinuje fyzickú
silu, techniku a~stratégiu s~cieľom získať kontrolu
nad súperom prostredníctvom povolených chmatov
a~akcií. Dvaja zápasníci sa snažia dosiahnuť víťazstvo
buď položením súpera na lopatky, alebo získaním
väčšieho počtu bodov za technické akcie v priebehu
zápasu \cite{zapasenie_pravidla}.

\subsection{Ciele a~priebeh zápasenia}

Cieľom zápasenia je dosiahnuť víťazstvo nad súperom
buď pádom (obe lopatky súpera súčasne na žinenke),
čo okamžite ukončuje zápas, alebo získaním väčšieho
počtu bodov za technické akcie
\cite{zapasenie_pravidla}.

Zápas sa skladá z~dvoch polčasov po 3~minúty
s~30-sekundovou prestávkou medzi nimi. Zápasníci
začínajú v~stoji a~môžu pokračovať v~parteri
(zápas na zemi), ak jeden z~nich získa kontrolu nad
súperom. Zápas môže byť ukončený predčasne technickou
prevahou -- ak rozdiel v~skóre dosiahne 8~bodov v~grécko-rímskom
alebo 10~bodov vo voľnom štýle. V~prípade remízy sa víťaz
určuje podľa kritérií ako vyšší počet technických bodov
alebo posledná bodovaná akcia \cite{zapasenie_pravidla}.

\subsection{Úloha rozhodcu v~zápasení}

Rozhodca zohráva kľúčovú úlohu pri zabezpečení
férového a~bezpečného priebehu zápasu. Dohliada na
dodržiavanie pravidiel, kontroluje správanie
zápasníkov a~komunikuje s~ostatnými členmi
rozhodcovského tímu -- bodovým rozhodcom a~predsedom
žinenky \cite{zapasenie_pravidla}.

Rozhodca je zodpovedný za signalizáciu začiatku
a~konca zápasu, prerušenie v~prípade zranenia alebo
konzultácie s~rozhodcovským tímom, ako aj za
sledovanie a~hodnotenie technických akcií. Jeho
rozhodnutia sú záväzné, v~sporných situáciách však
môže predseda žinenky rozhodnutie prehodnotiť
\cite{zapasenie_pravidla}.

\subsection{Hodnotenie zápasu rozhodcom}

Rozhodca hodnotí zápas prideľovaním bodov za
technické akcie. Za jednoduchý hod udeľuje 2~body,
za vysoký hod s~amplitúdou 4~body a~za vytlačenie
súpera zo žinenky 1~bod. Bodové hodnotenie rozhodca
signalizuje zdvihnutím ruky s~príslušným počtom
prstov \cite{zapasenie_pravidla}.

\subsection{Napomenutia a~zákazy}

Napomenutie je trest, ktorý rozhodca na žinenke udeľuje
zápasníkovi za porušenie pravidiel. V~bodovacích záznamoch
sa označuje symbolom \uv{O}. Samotné napomenutie neznamená
priamu stratu bodov pre previnilého zápasníka -- dôsledkom
napomenutia je však automatické udelenie bodov jeho súperovi.
Počet udelených bodov závisí od typu priestupku a~od štýlu
zápasenia \cite[čl.~54]{zapasenie_pravidla}.

Z~pohľadu rozhodcu na žinenke je napomenutie vždy spojené
s~následným bodovaním. Rozhodca najprv gestom signalizuje
napomenutie -- ukáže na previnilého zápasníka -- a~následne
pomocou gesta bodovanie oznámi, koľko bodov sa udeľuje
súperovi. Tieto dva signály tvoria dôležitú dvojicu,
ktorá je kľúčová aj pre návrh gest v~našej aplikácii.

Napomenutie sa podľa Medzinárodných pravidiel
zápasenia~\cite{zapasenie_pravidla} udeľuje za nasledujúce
typy priestupkov:

\begin{itemize}
  \item \textbf{Únik z~chmatu} -- brániaci zápasník
        sa otvorene vyhýba kontaktu so súperom, aby
        mu zabránil vykonať chmat
        \cite[čl.~50]{zapasenie_pravidla}.
  \item \textbf{Útek zo žinenky} -- zápasník opúšťa
        priestor žinenky z~postoja alebo z~partera
        \cite[čl.~51]{zapasenie_pravidla}.
  \item \textbf{Zakázané chmaty} -- napríklad držanie
        za hrdlo, vykrúcanie paže o~viac ako 90°,
        škrtenie alebo prelamovanie mostu tlakom
        k~hlave \cite[čl.~52]{zapasenie_pravidla}.
  \item \textbf{Všeobecné zákazy} -- napríklad ťahanie
        za vlasy, hryznutie, kopanie, udieranie do
        hlavy, či dohoda na výsledku zápasu
        \cite[čl.~49]{zapasenie_pravidla}.
\end{itemize}

Detailné podmienky a~výška udeľovaných bodov za jednotlivé
priestupky sú definované v~kapitole~9 Medzinárodných pravidiel
zápasenia~\cite{zapasenie_pravidla}.

\subsection{Pasivita}

Pasivita je situácia, keď sa zápasník úmyselne vyhýba
boju -- blokuje súpera, zaberá prsty alebo sa snaží
naťahovať čas. Postup penalizácie sa líši podľa štýlu
zápasenia. Vo voľnom štýle rozhodca po druhom
upozornení určí pasívnemu zápasníkovi 30-sekundovú
dobu aktivity. Ak počas nej neskóruje, jeho súper
získa 1~bod. V~grécko-rímskom štýle je postup
postupný -- prvé upozornenie je slovné, druhé sa
zapisuje ako prvé \uv{P} bez bodov a~tretie
previnenie je penalizované 1~bodom pre súpera
\cite[čl.~48]{zapasenie_pravidla}.

Rozhodca na žinenke signalizuje pasivitu ukázaním
ruky smerom k~pasívnemu zápasníkovi. Keďže je
potrebné rozlíšiť, ktorý zápasník je pasívny,
v~našej aplikácii sme pre toto gesto navrhli dve
alternatívy -- pre červeného a~modrého zápasníka.

\subsection{TOUCHE -- lopatkové víťazstvo}

TOUCHE, označované aj ako \uv{víťazstvo na lopatky},
je najvyššou formou víťazstva v~zápasení. Nastáva
vtedy, keď útočiaci zápasník drží súpera na žinenke
tak, že obe jeho lopatky sú súčasne v~kontakte so
žinenkou po určitý čas. Dosiahnutie TOUCHE okamžite
ukončuje zápas bez ohľadu na aktuálne skóre
\cite{zapasenie_pravidla}.

Podľa článku 44 Medzinárodných pravidiel zápasenia
\cite{zapasenie_pravidla} musí byť lopatkové
víťazstvo zreteľné -- obidve ramená súpera musia
byť súčasne pritlačené na žinenku v~pásme aktivity
alebo v~oranžovej zóne. Lopatky v~ochrannom pásme
sú neplatné.

Na rozdiel od bodovania alebo napomenutia, TOUCHE
vyžaduje potvrdenie od predsedu žinenky. Rozhodca
na žinenke najprv signalizuje TOUCHE, počká na
potvrdenie a~až potom udrie dlaňou po žinenke
a~odpíska koniec zápasu \cite{zapasenie_pravidla}.

\subsection{Analýza pravidiel z~pohľadu informatiky}
\label{subsec:analyza-informaticka}

Pravidlá popísané v~predchádzajúcich podsekciách
sme analyzovali z~pohľadu návrhu aplikácie na
inteligentné hodinky. Z~tejto analýzy vyplynuli
tri kľúčové pozorovania, ktoré priamo ovplyvnili
návrh nášho riešenia.

\subsubsection{Konečná množina sémanticky odlíšených udalostí}
Rozhodca počas zápasu opakovane vykonáva obmedzenú
sadu akcií -- bodovanie (1, 2 alebo 4 body pre
červeného alebo modrého zápasníka), napomenutie,
pasivita a~TOUCHE. Každá z~týchto udalostí má jasne
definovanú sémantiku v~kontexte stavu zápasu, čo
umožňuje modelovať úlohu ako klasifikačný problém
nad konečnou množinou tried.

\subsubsection{Jednoruké snímanie a rozlíšenie strán}
V~tradičnej signalizácii rozhodca rozlišuje strany
zápasníkov pomocou manžiet odlišnej farby na ľavom
(červená) a~pravom (modrá) zápästí. Naša aplikácia
však sníma pohyb len z~jedného zápästia, čím sa
vizuálne rozlíšenie strán stráca. Túto stratu
informácie kompenzujeme návrhom samostatných gest
pre každý farebný variant, pričom gestá, ktoré
rozhodca tradične vykonáva druhou rukou,
v~našej aplikácii simulujeme rozdielnou polohou
ľavej ruky -- typicky pred bruchom alebo za chrbtom
(viď kapitolu~\ref{chapter:gesta}).

\subsubsection{Kompozitnosť signálov a~závislosť na potvrdení}
Niektoré akcie rozhodcu nie sú samostatné, ale
predstavujú postupnosť gest -- napríklad napomenutie
je vždy nasledované bodovaním pre súpera. Iné akcie
(TOUCHE) sú zase platné len po potvrdení predsedom
žinenky. Gesto TOUCHE sa po detekcii zapisuje len
ako informačný záznam v~aplikácii, keďže o~platnosti
lopatkového víťazstva rozhoduje predseda žinenky
mimo nášho systému.